%iffalse
\let\negmedspace\undefined
\let\negthickspace\undefined
\documentclass[journal,12pt,twocolumn]{IEEEtran}
\usepackage{cite}
\usepackage{amsmath,amssymb,amsfonts,amsthm}
\usepackage{algorithmic}
\usepackage{graphicx}
\usepackage{textcomp}
\usepackage{xcolor}
\usepackage{txfonts}
\usepackage{listings}
\usepackage{enumitem}
\usepackage{mathtools}
\usepackage{gensymb}
\usepackage{comment}
\usepackage[breaklinks=true]{hyperref}
\usepackage{tkz-euclide} 
\usepackage{listings}
\usepackage{gvv}                                        
%\def\inputGnumericTable{}                                 
\usepackage[latin1]{inputenc}                                
\usepackage{color}                                            
\usepackage{array}                                            
\usepackage{longtable}                                       
\usepackage{calc}                                             
\usepackage{multirow}                                         
\usepackage{hhline}                                           
\usepackage{ifthen}                                           
\usepackage{lscape}
\usepackage{tabularx}
\usepackage{array}
\usepackage{float}


\newtheorem{theorem}{Theorem}[section]
\newtheorem{problem}{Problem}
\newtheorem{proposition}{Proposition}[section]
\newtheorem{lemma}{Lemma}[section]
\newtheorem{corollary}[theorem]{Corollary}
\newtheorem{example}{Example}[section]
\newtheorem{definition}[problem]{Definition}
\newcommand{\BEQA}{\begin{eqnarray}}
\newcommand{\EEQA}{\end{eqnarray}}
\newcommand{\define}{\stackrel{\triangle}{=}}
\theoremstyle{remark}
\newtheorem{rem}{Remark}

% Marks the beginning of the document
\begin{document}
\bibliographystyle{IEEEtran}
\vspace{3cm}

\title{Chapter 4 Permutation and Combination}
\author{ee24btech11012 - Bhavanisankar}
\maketitle
\newpage
\bigskip

\renewcommand{\thefigure}{\theenumi}
\renewcommand{\thetable}{\theenumi}

 6 A student is to answer 10 out of 13 questions in an examination such that he must choose atleast four from the first five questions. The number of choices available to him is  \hfill[2003]
\begin{enumerate} 
    \item  346
    \item  140 
    \item  196 
    \item  280
\end{enumerate}
7.The number of ways in which 6 men and 5 women can dine around a round table if no two women are to sit together is given by \hfill[2003]
\begin{enumerate}
 \item $7!\times5!$
 \item $6!\times 5!$
 \item $30!$
 \item $5!\times4!$
 \end{enumerate}
 8. How many ways are there to arrange the letters in the word GARDEN with vowels iin alphabetical order \hfill[2004]
 \begin{enumerate}
     \item 480
     \item 240
     \item 360
     \item 120
 \end{enumerate}
 9. The number of ways of distributing 8 identical balls in 3 distinct boxes so that none of the boxes is empty is \hfill[2004]
 \begin{enumerate}
     \item $ 8\choose3$
     \item 21
     \item $3^8$
     \item 5
 \end{enumerate}
 10. If the letters of the word SACHIN are arranged in all the possible ways and these words are written out as in dictionary, then the word SACHIN appears at the serial number \hfill[2005]
 \begin{enumerate}
     \item 601
     \item 600
     \item 603
     \item 602
 \end{enumerate}
11. At an election, a voter may vote for any number of candidates, not greater than the number to be elected. There are 10 candidates and 4 are of be selected, if a voter votes for atleast one candidate, then the number of ways in which he can vote is \hfill[2006]
\begin{enumerate}
 \item 5040
 \item 6210
 \item 385
 \item 1110
 \end{enumerate}
 12. The set $S={1,2,3,...,12}$ is to be partitioned into three sets A,B and C of equal size. Thus, $A\bigcup B\bigcup C=S$, $A \bigcap B=B \bigcap  C =A \bigcap C=\phi$. The number of ways to partition S is \hfill[2007]
 \begin{enumerate}
     \item $ \frac{12!}{(4!)^3}$
     \item $\frac{12!}{(4!)^4}$
     \item $\frac{12!}{3!(4!)^3}$
     \item $\frac{12!}{3!(4!)^4}$
 \end{enumerate}
 13. In a shop, there are five types of ice-creams available.A child buys six ice-creams.\\\textbf{Statement1}: The number of different ways in which the child can buy six ice-creams is $10\choose5$.\\\textbf{Statement2}:The number of different ways in which the child can buy six ice-creams is equal to the number of different ways of arranging 6 A's and 4 B's in a row. \hfill[2008]
 \begin{enumerate}
     \item Statement 1 is false, Statement 2 is true.
     \item Statement 1 is true, Statement 2 is true, Statement 2 is the correct explanation of Statement 1.
     \item Statement 1 is true, Statement 2 is true, Statement 2 is not a correct explanation of statement 1.
     \item Statement 1 is true, Statement 2 is false.
 \end{enumerate}
 14.How many different words can be formed by jumbling the letters in the word MISSISSIPPI in which no two S are adjacent ?
 \begin{enumerate}
     \item 8$6\choose4$$7\choose4$
     \item $6\times7$$8\choose4$
     \item $6\times8$$7\choose4$
     \item 7$6\choose4$$8\choose4$
 \end{enumerate}
 
\end{document}
